\scanpage{39}
Donc $U_\rho$ opère \inserted{\unsure{fidèlt|fidèlem.|fidèlement}} sur $\unsure{E|\mathcal E},E'$ par

\[
\left\{\begin{aligned}
u_\lambda^E &=\mathrm{id}+\lambda\pi'\otimes\pi\\
u_\lambda^{E'}&=\mathrm{id}+\lambda\pi'\otimes\pi
\end{aligned}\right.
\]

(\unsure{représentations|opérations} \deleted{\unsure{contragrédientes}}\inserted{transposées} l'une de l'autre) Comme $\pi,\pi'$ \unsure{invariants|invariantes} sous les $\tau_i$, \deleted{\unsure{donc pour|donc aussi}} sous les $\varphi^E(s)$, $\varphi_{\unsure{*|s}}^{E'}(s)$ $(s\in G_n)$, on trouve

\[
\begin{aligned}
\unsure{\alpha_2|\alpha'|(*)}\quad U_\rho&\longrightarrow\underline{\operatorname{Aut}}(E,\varphi^E)\\
\alpha\quad U_\rho&\longrightarrow\underline{\operatorname{Aut}}(X,\varphi^X)
\end{aligned}
\]

et de $=$ avec $E'$. On obtient en fait, en utilisant les scalaires

\diagram{\begin{tikzcd}[ampersand replacement=\&,column sep=small,row sep=small]
(\unsure{\alpha\alpha|\alpha *|**}) \& G_m\times U_\rho\arrow[rr]\&\&\operatorname{Aut}(E,\varphi^E)\\
\&\&G_m\arrow[ul,hook']\arrow[ur,swap,"\text{opère par homothéties}"]\&
\end{tikzcd}}

dont $\alpha$ est déduit en passant au quotient. On trouve

\uline{Th\unsure{m|me}} $\alpha$ et $(\alpha\alpha)$ sont des iso.

De plus, si $D_0=(H_1\cap\cdots\cap H_n)$, $D_0^{**}=D_0\setminus\{c,d\}$, où $\unsure{\exists|\,|\{}d=D_0\cap C$ (qui est une section de $X$ sur $S$) \unsure{.|:} on a que
% archive sheet number: 38
